\documentclass[11pt,a4paper]{article}

\usepackage[T1]{fontenc}
\usepackage{lmodern}
\usepackage[a4paper,margin=1in]{geometry}
\usepackage{amsmath,amssymb,amsthm}
\usepackage{enumitem}

\newtheorem{definition}{Definition}

\newcommand{\Grammar}{\mathcal{G}}
\newcommand{\NonTerminals}{\mathcal{N}}
\newcommand{\Terminals}{\Sigma}
\newcommand{\Rules}{\mathcal{P}}
\newcommand{\Items}{\mathcal{I}}
\newcommand{\Edits}{\mathcal{O}}

\title{A Formal Definition of ALCIEP}
\author{}
\date{}

\begin{document}
\maketitle

Let
\[
  \Grammar=(\NonTerminals,\Terminals,\Rules,S)
\]
be a context-free grammar and let
$w=t_1\cdots t_n\in\Terminals^{*}$ be the input token sequence.

An Earley item has the form
\[
  [A\to\alpha\,\bullet\,\beta, k],
\]
where $A\to\alpha\beta\in\Rules$ and $k$ is the index of the Earley set in
which the item was introduced.  The item is complete if $\beta=\varepsilon$;
otherwise, the symbol immediately after the dot is its expectation.

An Earley set $E_i\subseteq\Items$ is the set of items reachable after the
parser has performed $i$ correction operations.  The input position is kept
separately as $j$, the index of the next unread input token.  Let
\[
  Q_i=\{A\in\Terminals\mid
       [X\to\alpha\,\bullet\,A\beta,j]\in E_i\}
\]
be the set of terminals expected by $E_i$.

\begin{definition}[All-Correction Interactive Earley Parser]
  The All-Correction Interactive Earley Parser (ALCIEP) is the transition
  system whose configurations are tuples
  \[
    \langle E_i,j,\gamma\rangle,
  \]
  where $E_i$ is the current Earley set, $j$ is the index of the next unread
  input token, and $\gamma\in\Edits^{*}$ is the sequence of operations
  selected so far.  The Earley set is closed under the following two rules,
  which do not consume input and do not append an edit operation:
  \begin{itemize}[leftmargin=*]
    \item \textbf{Predictor:} If
    $[A\to\alpha\,\bullet\,B\beta,j]\in E_i$ and
    $B\to\delta\in\Rules$ for $B\in\NonTerminals$, then add
    \[
      [B\to\bullet\,\delta,i]\in E_i.
    \]

    \item \textbf{Completion:} If
    $[B\to\delta\,\bullet,k]\in E_i$ and
    $[A\to\alpha\,\bullet\,B\beta,j]\in E_k$, then add
    \[
      [A\to\alpha B\,\bullet\,\beta,j]\in E_i.
    \]
  \end{itemize}
  After this closure, the admissible interactive transitions are:
  \begin{itemize}[leftmargin=*]
    \item \textbf{Read:} If $j<n$ and $t_{j+1}$ matches $A\in Q_i$, advance
    the dot over $A$ and consume $t_{j+1}$:
    \[
      \langle E_i,j,\gamma\rangle
      \xrightarrow{\mathsf{read}(t_{j+1})}
      \langle E_{i+1},j+1,\gamma\mathsf{read}(t_{j+1})\rangle.
    \]

    \item \textbf{Replacement:} If $j<n$ and $t_{j+1}$ does not match
    $A\in Q_i$, advance the dot over $A$ and consume $t_{j+1}$:
    \[
      \langle E_i,j,\gamma\rangle
      \xrightarrow{\mathsf{replace}(t_{j+1},A)}
      \langle E_{i+1},j+1,
        \gamma\mathsf{replace}(t_{j+1},A)\rangle.
    \]

    \item \textbf{Deletion:} If $j<n$, consume $t_{j+1}$ without advancing
    any Earley item:
    \[
      \langle E_i,j,\gamma\rangle
      \xrightarrow{\mathsf{delete}(t_{j+1})}
      \langle E_i,j+1,\gamma\mathsf{delete}(t_{j+1})\rangle.
    \]

    \item \textbf{Insertion:} For every $A\in Q_i$, advance the dot over
    $A$ without consuming an input token:
    \[
      \langle E_i,j,\gamma\rangle
      \xrightarrow{\mathsf{insert}(A)}
      \langle E_{i+1},j,\gamma\mathsf{insert}(A)\rangle.
    \]
  \end{itemize}
  A configuration is accepting if $j=n$ and $E_i$ contains a complete start
  item $[S\to\alpha\,\bullet,0]$.  The sequence $\gamma$ of selected edit
  operations is then the correction produced by ALCIEP.
\end{definition}

\end{document}
